\documentclass{article}
\usepackage[utf8]{inputenc}
\usepackage{amsmath,amssymb}
\usepackage[most]{tcolorbox}
\usepackage{geometry}
\geometry{margin=2.5cm}

\newcommand{\step}[1]{\par\noindent\hangindent=3.5em\hangafter=1%
  \makebox[3.5em][l]{\textbf{#1.}}}
\newcommand{\steptag}[1]{\hfill{\small\textsf{[#1]}}}

\newtcolorbox{proofbox}[1]{
  colback=gray!5, colframe=gray!60,
  fonttitle=\bfseries, title={#1},
  breakable, enhanced,
  left=4pt, right=4pt, top=4pt, bottom=4pt,
  before skip=10pt, after skip=10pt
}

\begin{document}

\begin{proofbox}{Group-input polar closeness: there exist universal C\_grp,...}
\step{1} Group-input polar closeness: there exist universal C\_grp, C\_pol >= 1, kappa\_pol in (0, 1/2] such that for every finite-dimensional exact-unit epsilon\_r-C*-algebra and delta > 0 satisfying C\_pol*(epsilon\_r + delta) <= kappa\_pol and C\_grp*epsilon\_r < delta - C\_pol*(epsilon\_r*delta + delta\textasciicircum{}2), the inverse u\_delta:S\_delta -> calU of the polar map satisfies ||u\_delta(U bold-dot V) - U bold-dot V|| <= C\_grp*epsilon\_r and ||u\_delta(U\textasciicircum{}dagger) - U\textasciicircum{}dagger|| <= C\_grp*epsilon\_r for every U, V in calU. \steptag{validated} \\[6pt]
\step{1.1} Constant choice and polar range: take C\_pol\textasciicircum{}0,kappa\_pol\textasciicircum{}0 from the validated external lem-stage1-polar-retraction, set C\_pol=C\_pol\textasciicircum{}0, kappa\_pol=min(kappa\_pol\textasciicircum{}0,1/16), C\_grp=5, and rho=delta-C\_pol*(epsilon\_r*delta+delta\textasciicircum{}2). These are universal with C\_grp,C\_pol>=1 and kappa\_pol in (0,1/2]. Under the root hypotheses, epsilon\_r+delta<=1/16 and rho>5*epsilon\_r>=0; lem-stage1-polar-retraction applies and gives calU\_rho subseteq S\_delta together with X=u\_delta(X) bold-dot h\_delta(X), u\_delta(X) in calU, and ||h\_delta(X)-J||<delta for X in S\_delta. \steptag{validated} \\[4pt]
\step{1.2} Raw input estimates and domain typing: with the constants and rho of 1.1, every U,V in calU satisfy ||(U bold-dot V)\textasciicircum{}dagger bold-dot (U bold-dot V)-J||<=3*epsilon\_r and U bold-dot V has a right inverse, while ||(U\textasciicircum{}dagger)\textasciicircum{}dagger bold-dot U\textasciicircum{}dagger-J||<=2*epsilon\_r and U\textasciicircum{}dagger has the right inverse U. Consequently U bold-dot V and U\textasciicircum{}dagger both belong to calU\_rho and hence to S\_delta. \steptag{validated} \\[6pt]
\step{1.2.1} Unitary norm and right-inverse data: for W in calU, def-approximate-unitary-space gives W\textasciicircum{}dagger bold-dot W=J and an R with W bold-dot R=J. The exact-unit clause of def-epsilon-cstar-algebra gives ||J||=1, while its C*-lower bound gives 1=||W\textasciicircum{}dagger bold-dot W||>=(1-epsilon\_r)||W||\textasciicircum{}2; since epsilon\_r<=1/16, ||W||\textasciicircum{}2<=1/(1-epsilon\_r)<=16/15. \steptag{validated} \\[4pt]
\step{1.2.2} Opposite-product defect: for W in calU and R as in 1.2.1, the associator axiom gives ||R-W\textasciicircum{}dagger||=||(W\textasciicircum{}dagger bold-dot W) bold-dot R-W\textasciicircum{}dagger bold-dot (W bold-dot R)||<=epsilon\_r||W||\textasciicircum{}2||R||. Thus epsilon\_r||W||\textasciicircum{}2<=1/15 implies ||R||<=15||W||/14. Using J=W bold-dot R, the product-norm axiom, and ||W\textasciicircum{}dagger||=||W||, one obtains ||W bold-dot W\textasciicircum{}dagger-J||<=epsilon\_r(1+epsilon\_r)||W||\textasciicircum{}3||R||<=((17/16)*(15/14)*(256/225))*epsilon\_r<=2*epsilon\_r. \steptag{validated} \\[4pt]
\step{1.2.3} Product defect: for U,V in calU and X=U bold-dot V, the involution axiom gives X\textasciicircum{}dagger=V\textasciicircum{}dagger bold-dot U\textasciicircum{}dagger. Two applications of the associator axiom compare (V\textasciicircum{}dagger bold-dot U\textasciicircum{}dagger) bold-dot (U bold-dot V) first with ((V\textasciicircum{}dagger bold-dot U\textasciicircum{}dagger) bold-dot U) bold-dot V and then with (V\textasciicircum{}dagger bold-dot (U\textasciicircum{}dagger bold-dot U)) bold-dot V=J. The product-norm axiom and 1.2.1 therefore give ||X\textasciicircum{}dagger bold-dot X-J||<=2*epsilon\_r*(1+epsilon\_r)||U||\textasciicircum{}2||V||\textasciicircum{}2<=(544/225)*epsilon\_r<=3*epsilon\_r. \steptag{validated} \\[4pt]
\step{1.2.4} Exact right inverse for the product: for X=U bold-dot V, the product-norm estimate and 1.2.1 give ||X||\textasciicircum{}2<=(1+epsilon\_r)\textasciicircum{}2*(16/15)\textasciicircum{}2<=289/225. If X bold-dot Y=0, the associator axiom and the defect bound of 1.2.3 imply ||Y||<=||(J-X\textasciicircum{}dagger bold-dot X) bold-dot Y||+||(X\textasciicircum{}dagger bold-dot X) bold-dot Y-X\textasciicircum{}dagger bold-dot (X bold-dot Y)||<=((1+epsilon\_r)*3*epsilon\_r+epsilon\_r||X||\textasciicircum{}2)||Y||<=(51/256+289/3600)||Y||<(1/2)||Y||. Hence Y=0. Bilinearity makes the left multiplier L\_X linear; on the finite-dimensional algebra it is therefore surjective, so some R\_X satisfies X bold-dot R\_X=J. \steptag{validated} \\[4pt]
\step{1.2.5} Domain conclusion: X\_*=U\textasciicircum{}dagger has the exact right inverse U, and by the involution axiom its defect is ||U bold-dot U\textasciicircum{}dagger-J||<=2*epsilon\_r from 1.2.2. The product X\_p=U bold-dot V has a right inverse by 1.2.4 and defect at most 3*epsilon\_r by 1.2.3. Since rho>5*epsilon\_r, each defect is strictly below 2*rho (if epsilon\_r=0, use rho>0 separately). By def-approximate-unitary-space both inputs lie in calU\_rho, and 1.1 (lem-stage1-polar-retraction) gives calU\_rho subseteq S\_delta. \steptag{validated} \\[4pt]
\step{1.3} Polar closeness estimate: under epsilon\_r+delta<=1/16, if X lies in S\_delta and ||X\textasciicircum{}dagger bold-dot X-J||<=3*epsilon\_r, then the polar inverse from lem-stage1-polar-retraction satisfies ||u\_delta(X)-X||<=5*epsilon\_r. \steptag{validated} \\[6pt]
\step{1.3.1} Polar square comparison: write u=u\_delta(X), h=h\_delta(X), and a=h-J. By lem-stage1-polar-retraction, X=u bold-dot h, u lies in calU, h is Hermitian, and ||a||<delta. Hence ||u||\textasciicircum{}2<=1/(1-epsilon\_r)<=16/15 by the C*-lower bound applied to u\textasciicircum{}dagger bold-dot u=J, while ||h||<=1+delta<=17/16. The involution axiom gives X\textasciicircum{}dagger=h bold-dot u\textasciicircum{}dagger. Comparing (h bold-dot u\textasciicircum{}dagger) bold-dot (u bold-dot h) with ((h bold-dot u\textasciicircum{}dagger) bold-dot u) bold-dot h and then with (h bold-dot (u\textasciicircum{}dagger bold-dot u)) bold-dot h=h bold-dot h, the associator and product-norm axioms give ||X\textasciicircum{}dagger bold-dot X-h bold-dot h||<=2*epsilon\_r*(1+epsilon\_r)||h||\textasciicircum{}2||u||\textasciicircum{}2<=(9826/3840)*epsilon\_r<=3*epsilon\_r. \steptag{validated} \\[4pt]
\step{1.3.2} Quadratic absorption: exact unitality and bilinearity give h bold-dot h-J=(J+a) bold-dot (J+a)-J=2a+a bold-dot a. Combining the assumed defect ||X\textasciicircum{}dagger bold-dot X-J||<=3*epsilon\_r with 1.3.1 yields 2||a||<=6*epsilon\_r+(1+epsilon\_r)||a||\textasciicircum{}2. Since ||a||<delta<=1/16 and epsilon\_r<=1/16, (1+epsilon\_r)||a||\textasciicircum{}2<=(17/256)||a||. Thus (495/256)||a||<=6*epsilon\_r, so ||a||<=(1536/495)*epsilon\_r<=4*epsilon\_r. \steptag{validated} \\[4pt]
\step{1.3.3} Return to the polar factor: exact right unitality and X=u bold-dot h give ||u\_delta(X)-X||=||u bold-dot J-u bold-dot h||=||u bold-dot (J-h)||<=(1+epsilon\_r)||u||||a||. Using ||u||<=sqrt(16/15), epsilon\_r<=1/16, and 1.3.2 gives ||u\_delta(X)-X||<=(17/sqrt(15))*epsilon\_r<=5*epsilon\_r. \steptag{validated} \\[4pt]
\end{proofbox}

\end{document}
